% Sorgente LaTeX canonico, importato e revisionato dal precedente sorgente Markdown.
\chapter{Taylor e sviluppi di Maclaurin}
\label{chap:12}

\begin{obiettivocapitolo}{scegliere centro e ordine di Taylor per limiti e approssimazioni}
\prioritaalta. Taylor è un metodo di calcolo quando occorre controllare cancellazioni o quantificare l'errore.
\end{obiettivocapitolo}
\begin{proceduraesame}{progettare lo sviluppo}
\begin{enumerate}
\item Porta il punto di sviluppo a $h=x-x_0$ e individua la quantità richiesta.
\item Stima in anticipo l'ordine minimo: dopo una divisione per $h^m$ servono termini almeno fino a $h^m$; in una cancellazione potrebbe servire oltre.
\item Sviluppa ogni funzione allo stesso ordine effettivo, includendo composizioni e prodotti.
\item Semplifica e arrestati al primo termine non nullo che decide il risultato.
\item Per un'approssimazione numerica, usa il resto appropriato e dimostra che è sotto la tolleranza.
\end{enumerate}
\end{proceduraesame}
\begin{sceltametodo}{forma del resto}
Peano è adatto ai limiti locali; Lagrange e Cauchy alle stime puntuali; il resto integrale è utile quando si controlla bene la derivata successiva su un intervallo.
\end{sceltametodo}
\begin{errorefrequente}{serie di Taylor e funzione}
Essere $C^\infty$ non garantisce che la serie di Taylor converga alla funzione. Nei problemi locali finiti basta il polinomio con resto; non occorre assumere analiticità.
\end{errorefrequente}
\begin{controllofinale}{approssimazione}
Controlla segno, ordine del resto, numero di cifre affidabili e distanza dal centro di sviluppo.
\end{controllofinale}

\(T_{n,x_0}f\) è il polinomio di Taylor di grado al più \(n\) centrato in \(x_0\); Maclaurin è il caso particolare \(x_0=0\). Il termine \(R_n=f-T_{n,x_0}f\) è il resto.

\section{Polinomio e coefficienti di Taylor}\label{polinomio-e-coefficienti-di-taylor}

\[
\displaystyle
T_{n,x_0}f(x)
=
\sum_{k=0}^{n}
\frac{f^{(k)}(x_0)}{k!}(x-x_0)^k.
\]

Se

\[
\displaystyle
T_{n,x_0}f(x)=\sum_{k=0}^{n}c_k(x-x_0)^k,
\]

allora

\[
\displaystyle
c_k=\frac{f^{(k)}(x_0)}{k!},
\qquad
f^{(k)}(x_0)=k!\,c_k.
\]

Approfondimento: \href{https://ingegnerismo.it/matematica/formula-di-taylor/}{formula di Taylor}.

\section{Resti di Peano, Lagrange, Cauchy e integrale}\label{resti-di-peano-lagrange-cauchy-e-integrale}

Una condizione sufficiente per la forma di Peano è \(f\in C^n\) in un intervallo aperto contenente \(x_0\). Allora, per \(x\to x_0\):

\[
\displaystyle
f(x)=T_{n,x_0}f(x)+o((x-x_0)^n).
\]

Per \(x\ne x_0\), se \(f,f',\ldots,f^{(n)}\) sono continue sul segmento di estremi \(x_0,x\) e \(f^{(n+1)}\) esiste al suo interno, esiste \(\xi_L\) strettamente interno al segmento tale che

\[
\displaystyle
f(x)=T_{n,x_0}f(x)
+\frac{f^{(n+1)}(\xi_L)}{(n+1)!}(x-x_0)^{n+1}.
\]

Se \(|f^{(n+1)}(t)|\le M\) sul segmento:

\[
\displaystyle
|f(x)-T_{n,x_0}f(x)|
\le
\frac{M}{(n+1)!}|x-x_0|^{n+1}.
\]

Sotto le stesse ipotesi, esiste un punto \(\xi_C\) strettamente interno al segmento tale che:

\[
R_n(x)
=
\frac{f^{(n+1)}(\xi_C)}{n!}
(x-\xi_C)^n(x-x_0).
\]

Se \(f\in C^{n+1}\) sul segmento di estremi \(x_0\) e \(x\):

\[
R_n(x)
=
\frac1{n!}\int_{x_0}^{x}f^{(n+1)}(t)(x-t)^n\,dt,
\]

\[
f(x)=T_{n,x_0}f(x)+R_n(x).
\]

La forma integrale richiede integrale definito e teorema fondamentale del calcolo. Approfondimento: \href{https://ingegnerismo.it/matematica/resto-di-taylor/}{resto di Taylor}. Esercizi: \href{https://ingegnerismo.it/matematica/resti-di-taylor-esercizi/}{resti di Taylor}.

\section{Progetto dell'approssimazione}\label{progetto-dellapprossimazione}

Se \(A=T_{n,x_0}f(x)\) e \(|f^{(n+1)}(t)|\le M_{n+1}\) sul segmento fra \(x_0\) e \(x\), porre

\[
B_n=\dfrac{M_{n+1}}{(n+1)!}|x-x_0|^{n+1}.
\]

Allora \(|f(x)-A|\le B_n\). Per una tolleranza assoluta \(\varepsilon\) si cerca il più piccolo \(n\) per cui \(B_n\le\varepsilon\); il fallimento di questa maggiorazione non prova, da solo, che l'errore effettivo superi la tolleranza. Se \(B_n<|A|\):

\[
\dfrac{|f(x)-A|}{|f(x)|}
\le
\dfrac{B_n}{|A|-B_n}.
\]

L'intervallo \([A-B_n,A+B_n]\) certifica le cifre che si arrotondano nello stesso modo in tutti i suoi punti; segno del resto e stime unilaterali possono restringerlo. Esercizi: \href{https://ingegnerismo.it/matematica/approssimazione-numerica-con-taylor-esercizi/}{approssimazione numerica con Taylor}.

\section{Sviluppi di Maclaurin notevoli}\label{sviluppi-di-maclaurin-notevoli}

Per \(x\to0\):

\[
\displaystyle
e^x=\sum_{k=0}^{n}\frac{x^k}{k!}+o(x^n),
\qquad n\ge0.
\]

\[
\displaystyle
\ln(1+x)
=
\sum_{k=1}^{n}(-1)^{k+1}\frac{x^k}{k}+o(x^n),
\qquad n\ge1,
\quad x>-1.
\]

Per \(\alpha\in\mathbb R\), \(n\ge0\) e \(1+x>0\):

\[
\displaystyle
(1+x)^\alpha
=
\sum_{k=0}^{n}\binom{\alpha}{k}x^k+o(x^n),
\]

\[
\displaystyle
\binom{\alpha}{0}=1,
\qquad
\binom{\alpha}{k}
=
\frac{\alpha(\alpha-1)\cdots(\alpha-k+1)}{k!}.
\]

Per ogni \(n\ge0\):

\[
\displaystyle
\sin x
=
\sum_{k=0}^{n}(-1)^k\frac{x^{2k+1}}{(2k+1)!}
+o(x^{2n+1}).
\]

\[
\displaystyle
\cos x
=
\sum_{k=0}^{n}(-1)^k\frac{x^{2k}}{(2k)!}
+o(x^{2n}).
\]

\[
\displaystyle
\tan x=x+\frac{x^3}{3}+\frac{2x^5}{15}+o(x^5).
\]

\[
\displaystyle
\arctan x=x-\frac{x^3}{3}+\frac{x^5}{5}+o(x^5).
\]

\[
\displaystyle
\arcsin x=x+\frac{x^3}{6}+\frac{3x^5}{40}+o(x^5).
\]

\[
\displaystyle
\arccos x
=\frac\pi2-x-\frac{x^3}{6}-\frac{3x^5}{40}+o(x^5).
\]

\[
\displaystyle
\sinh x
=
\sum_{k=0}^{n}\frac{x^{2k+1}}{(2k+1)!}
+o(x^{2n+1}).
\]

\[
\displaystyle
\cosh x
=
\sum_{k=0}^{n}\frac{x^{2k}}{(2k)!}
+o(x^{2n}).
\]

Approfondimento: \href{https://ingegnerismo.it/matematica/sviluppo-asintotico/}{sviluppo asintotico}.

\section{Sviluppi aggiuntivi di uso frequente}\label{sviluppi-aggiuntivi-di-uso-frequente}

Altri sviluppi utili, per \(x\to0\):

\[
\frac1{1-x}=\sum_{k=0}^{n}x^k+o(x^n),
\qquad
\frac1{1+x}=\sum_{k=0}^{n}(-1)^kx^k+o(x^n),
\]

\[
\ln(1-x)=-\sum_{k=1}^{n}\frac{x^k}{k}+o(x^n),
\]

\[
\sqrt{1+x}
=1+\frac x2-\frac{x^2}{8}+\frac{x^3}{16}+o(x^3),
\]

\[
\frac1{\sqrt{1+x}}
=1-\frac x2+\frac{3x^2}{8}-\frac{5x^3}{16}+o(x^3).
\]

\section{Serie di Taylor e analiticità}\label{serie-di-taylor-e-analiticituxe0}

Una funzione è \textbf{analitica reale in \(x_0\)} se esiste \(r>0\) e una serie di potenze convergente tale che

\[
f(x)=\sum_{k=0}^{\infty}c_k(x-x_0)^k
\qquad\text{per ogni }|x-x_0|<r.
\]

In tal caso la funzione è di classe \(C^\infty\) e necessariamente

\[
c_k=\frac{f^{(k)}(x_0)}{k!},
\]

perciò la serie è la serie di Taylor

\[
\displaystyle
\sum_{k=0}^{\infty}
\frac{f^{(k)}(x_0)}{k!}(x-x_0)^k.
\]

Per un punto \(x\) fissato nel dominio della serie di Taylor:

\[
\displaystyle
f(x)=\sum_{k=0}^{\infty}
\frac{f^{(k)}(x_0)}{k!}(x-x_0)^k
\quad\Longleftrightarrow\quad
R_n(x)=f(x)-T_{n,x_0}f(x)\to0
\quad(n\to\infty).
\]

\[
f\in C^\infty
\not\Longrightarrow
f(x)=\sum_{k=0}^{\infty}\frac{f^{(k)}(x_0)}{k!}(x-x_0)^k.
\]

Caso \(C^\infty\) non analitico in \(0\):

\[
\displaystyle
f(x)=
\begin{cases}
e^{-1/x^2}, & x\ne0,\\
0, & x=0,
\end{cases}
\qquad
f^{(k)}(0)=0\quad\forall k\ge0.
\]

Approfondimento: \href{https://ingegnerismo.it/matematica/serie-di-taylor/}{serie di Taylor}.
